\section{Introduction}
\label{sec:intro}

A legged robot's usefulness outside a laboratory rests on two things it cannot observe directly: whether its actuators
still deliver what the controller asks, and whether its feet are where the estimator believes they are. Both degrade
silently. An actuator loses torque gradually; a foot slips for a fraction of a step on a wet tile; an encoder drifts by
a few milliradians. None of these announce themselves, and each corrupts the state estimate and the controller that
depends on it.

% Fig. 1 — the whole claim in one single-column figure (Sprint 10 W2.2).
% Left: three structural invariances a healthy robot must satisfy. Middle: breaking one IS the fault/slip.
% Right: the three channels that read the break, and what each delivers.
\begin{figure}[t]
\centering
\begin{tikzpicture}[
  font=\scriptsize,
  box/.style={draw,rounded corners=2pt,align=center,inner sep=3pt,minimum height=8mm,text width=20mm},
  inv/.style={box,fill=blue!8,draw=blue!55},
  brk/.style={box,fill=red!8,draw=red!55,text width=20mm},
  ch/.style={box,fill=green!8,draw=green!50!black,text width=22mm},
  del/.style={align=left,inner sep=1pt,font=\tiny},
  ar/.style={-{Latex[length=1.6mm]},draw=black!60},
  lbl/.style={font=\tiny\bfseries,align=center}
]

% ---------------- column headings
\node[lbl] (h1) at (0,2.35) {structural invariance\\(healthy robot)};
\node[lbl] (h2) at (2.85,2.35) {breaking it\\\emph{is} the fault};
\node[lbl] (h3) at (6.05,2.35) {channel that reads it\\$\rightarrow$ what it delivers};

% ---------------- left: three invariances
\node[inv] (i1) at (0,1.45) {morphological\\$C_2$: $\mathrm{LF}\!\leftrightarrow\!\mathrm{RF}$, $\mathrm{LH}\!\leftrightarrow\!\mathrm{RH}$};
\node[inv] (i2) at (0,0.35) {gait space--time\\$\symgrp=\{e,(\gs,\tfrac{T}{2})\}$};
\node[inv] (i3) at (0,-0.85) {contact geometry\\rolling: $\dot d_i = R u_i$};

% tiny quadruped pictogram inside i1's margin
\begin{scope}[shift={(-1.25,1.45)},scale=0.16]
  \draw[black!55,fill=black!5] (-1.6,-0.7) rectangle (1.6,0.7);
  \draw[dashed,red!60] (0,-1.15) -- (0,1.15);
  \foreach \x/\y in {-1.15/0.75,1.15/0.75,-1.15/-0.75,1.15/-0.75} \draw[fill=blue!30] (\x,\y) circle (0.3);
\end{scope}

% ---------------- middle: the break
\node[brk] (b1) at (2.85,1.45) {one-sided fault\\(actuator, encoder,\\single-foot slip)};
\node[brk] (b2) at (2.85,0.35) {phase/duty break\\(gait asymmetry)};
\node[brk] (b3) at (2.85,-0.85) {constraint break\\(wheel slip, lost contact)};

% ---------------- right: channels
\node[ch] (c1) at (6.05,1.45) {$\Rminus$ on the raw element\\{\tiny(mirror flip test)}};
\node[ch] (c2) at (6.05,0.35) {$\Rminus$ on the \emph{residual}\\{\tiny(equivariant model)}};
\node[ch] (c3) at (6.05,-0.85) {invariant filter\\innovation / NIS};

% ---------------- arrows
\foreach \a/\b in {i1/b1,i2/b2,i3/b3} \draw[ar] (\a) -- (\b);
\foreach \a/\b in {b1/c1,b2/c2,b3/c3} \draw[ar] (\a) -- (\b);
% cross links: the residual channel also serves the one-sided fault; the filter also serves the gait break
\draw[ar,dashed,black!35] (b1) to[out=-25,in=170] (c2);
\draw[ar,dashed,black!35] (b3) to[out=25,in=190] (c2);

% ---------------- deliverables strip
\node[draw,fill=black!4,rounded corners=2pt,text width=82mm,align=left,inner sep=4pt,font=\tiny]
  (dl) at (3.0,-2.15) {%
  \textbf{What the construction delivers.}\;
  \emph{(i)} exact finite-sample false-alarm rate --- the test is a group-randomisation test, so its level holds at any
  $K$ with no fault data and no training (Thm.~\ref{thm:exact});\;
  \emph{(ii)} an anytime-valid sequential layer (e-process, Ville) that may be read at every cycle
  (Prop.~\ref{prop:seq});\;
  \emph{(iii)} a characterisation of \emph{what is visible}: a $\symgrp$-fixed (bilateral) fault is invisible to any
  invariance test, and the magnitude channel is what covers it (Thm.~\ref{thm:twolayer}).};

\begin{scope}[on background layer]
  \node[fit=(h1)(i3),draw=blue!25,rounded corners=3pt,inner sep=4pt,fill=blue!3] {};
  \node[fit=(h2)(b3),draw=red!25,rounded corners=3pt,inner sep=4pt,fill=red!3] {};
  \node[fit=(h3)(c3),draw=green!30!black!35,rounded corners=3pt,inner sep=4pt,fill=green!3] {};
\end{scope}
\end{tikzpicture}
\caption{\textbf{Fault and slip detection as invariance testing.} A healthy legged or wheeled-legged robot satisfies
three structural invariances (left): the morphological mirror symmetry of the body, the space--time symmetry of the
commanded gait, and the geometry of contact (a stance foot is fixed; a rolling wheel's contact point moves at the
measured rolling rate). Every fault or slip we target \emph{is} a break of one of them (middle), which turns detection
into a test of a symmetry hypothesis rather than a comparison against a learned or thresholded nominal. Three channels
read the break (right); dashed links mark the couplings that make them complementary rather than redundant.}
\label{fig:concept}
\end{figure}


The established answers fall into two families, and both leave the same hole. \emph{Model-based} residual methods run a
dynamics model --- a generalised-momentum observer is the canonical choice
\cite{deluca2003actuator,haddadin2017robot} --- and threshold the residual, usually against a $\chi^2$ quantile.
\emph{Data-driven} methods learn what nominal looks like, from an autoencoder or a recurrent classifier, and threshold a
reconstruction error or a predicted fault parameter. The hole is the false-alarm rate. A $\chi^2$ threshold is exact
only if the residual is i.i.d.\ Gaussian, and a legged robot's residual is neither: it is gait-periodic, serially
correlated, and non-stationary across a recording. We measure the consequence directly. On our simulation benchmark the
textbook momentum-observer $\chi^2$ detector runs at a nominal alarm rate of \textbf{0.15} against a nominal $0.05$, and
re-calibrating its threshold on the robot's own nominal data does not repair it --- on a \emph{held-out} stretch of the
same nominal recording the rate is \textbf{0.50}, because the residual's statistics have already moved
(\S\ref{sec:simulation}). On an external, third-party fault dataset a Mahalanobis detector fires on \textbf{52\,\%} of
healthy segments (0.522 false alarms per healthy gap) where our test fires on \textbf{4.9\,\%}, i.e.\ at its nominal
level (\S\ref{sec:external}). A detector whose false-alarm rate is unknown cannot be given authority over a controller,
which is why slip handling in practice is still a hand-tuned threshold on foot speed \cite{jenelten2019slippery}.

\paragraph{Gap}
Three literatures each hold one piece and none holds the combination. Morphological-symmetry work in robotics
\cite{ordonez2023discrete} uses a robot's discrete symmetry group to constrain models and policies, but treats symmetry
as an architectural prior, not as a hypothesis to be tested. Fault detection and contact estimation
\cite{camurri2017probabilistic,lin2021contact} build ever better nominal models and contact classifiers, but calibrate
their decisions empirically. Invariance testing in statistics \cite{hemerik2018exact,koning2023evalues} gives exactly
the guarantee the robotics side lacks --- a group-randomisation test whose level is exact in finite samples, and an
anytime-valid sequential form --- but has not been connected to a physical invariance a robot actually possesses. The
missing step is to notice that a healthy legged robot \emph{already satisfies} a group invariance precisely enough to
serve as the null hypothesis, and that the faults we care about are exactly the things that break it.

\paragraph{Thesis}
We propose to detect faults and slip by testing the structural invariances of the robot itself (Fig.~\ref{fig:concept}).
The morphological mirror symmetry of the body, the space--time symmetry of the commanded gait, and the geometry of
contact are properties a healthy machine must satisfy; a one-sided actuator fault, a single-foot slip, or a broken
rolling constraint are, by construction, violations of them. Framing detection this way replaces a learned or
thresholded nominal with a hypothesis that needs no fault data, no training, and --- in its model-free form --- no
dynamics model, and it inherits an exact finite-sample false-alarm rate from the group-randomisation construction.

\paragraph{Contributions}
\begin{enumerate}
\item \textbf{Invariance as the null, with an exact level.} We formulate the healthy-robot hypothesis $H_0$ as
  distributional invariance of a phase-registered data element under the gait symmetry group, and test it with a
  group-randomisation (flip) test whose false-alarm rate is \emph{exact at any sample size}, plus an anytime-valid
  sequential layer built from e-processes. We also give the recalibrated null $H_0'$ for the realistic case of a robot
  that is stably but measurably asymmetric, and a power lower bound with an explicit sample-size rule
  (\S\ref{sec:theory}; Thm.~\ref{thm:exact}, Prop.~\ref{prop:power}, Prop.~\ref{prop:seq}).
\item \textbf{A two-layer characterisation of what is visible, and the architecture it forces.} We prove a dichotomy:
  a fault fixed by the symmetry group (a bilaterally symmetric degradation) is invisible to \emph{every} invariance
  test, whatever its magnitude, provided the faulted loop keeps a unique symmetric attractor; anything that breaks the
  law is detectable by a consistent statistic. The mean-level layer gives the power of the specific paired statistic,
  and the two layers do not coincide --- a zero-mean law change is invisible to the paired statistic and visible to an
  energy-distance one. This is why the system needs three channels rather than one
  (\S\ref{sec:theory}; Thm.~\ref{thm:twolayer}).
\item \textbf{Equivariant nominal models, and why equivariance is not optional.} When a learned dynamics model is used
  to form a residual, we show the residual inherits $H_0$ only if the model is equivariant, quantify the contamination
  a non-equivariant model injects, and construct an equivariant Deep Lagrangian Network by mirror weight sharing. The
  effect is not marginal: a plain learned model drives the residual test's false-alarm rate to $1.00$ while its
  equivariant twin stays at $0.00$ (\S\ref{sec:theory},~\S\ref{sec:simulation}).
\item \textbf{Rolling contact preserves the invariant filter, and a five-layer evidence stack.} We prove that a rolling
  wheel's moving contact keeps the contact-aided invariant filter group-affine --- the invariant error dynamics is
  unchanged, only the mean and the process noise move --- which admits wheeled-legged robots to the same estimator and
  the same gating front-end as point-foot robots. The claims are evaluated at five levels of increasing externality,
  including \emph{two} of our own hardware platforms (\S\ref{sec:hardware}).
\end{enumerate}

\paragraph{What the evaluation shows}
On our own Unitree Go2 corpus --- eleven sessions, a machine over a year in service, three sites, two dates two months
apart --- the naive invariance test rejects on \textbf{11/11} sessions, i.e.\ it sees the machine's real, stable
left--right asymmetry, and the recalibrated level $\nu_0$ reproduces across the two-month gap to within a factor
\textbf{1.18}, which is smaller than the session-to-session spread within a single day. On the wheeled-legged M1 the
rolling-contact filter recovers \textbf{0.99--1.00} of the driven path length where a fixed-foot filter recovers
\textbf{0.03--0.04}. Used as an estimator front-end, the calibrated per-event gate discards \textbf{1.9\,\%} of nominal
contact measurements where the standard fixed foot-speed threshold discards \textbf{40.6\,\%} of them on the same real
outdoor data. And on the external benchmark the test holds its level at \textbf{0.049} against a nominal $0.05$.
